% v18. Exact exclusion of G_{5,5}; no denominator cutoff.
% Uses standard positive two-high form, no-negative-low-wedge,
% common-low reverse bound, and low-edge D opposition only at 23 and 24.

\section{五頂点六辺の house 型 support の排除}
\label{hs18:section}

三角形 $012$ と四角形 $0134$ が辺 $01$ を共有する graph
\[
 H=\{01,02,03,12,14,34\}=G_{5,5}
\]
を扱う．二価 cores は $2,3,4$ であり，従来の四頂点障害とは異なり，
全五頂点を使用する．外側の二価 cores $3,4$ の高隣接頂点の間は低辺なので，
補題~\ref{kb15:no-negative-low-wedge} により，この二つは正の枝に限られる．
中央の core $2$ は三角形上にあり，その負の枝をあらかじめ捨てない．

\subsection{共有高辺の共通低 ray の三つの型}

二高分母の canonical 標識を $(A,B,C,D)$ とし，
\[
 P=\begin{pmatrix}
 0&-A_0&x&-s\\
 -B_0&0&-y&z\\
 r&-r-1&0&0\\
 -S&T&-L&0
 \end{pmatrix},\qquad
 L=T-S+1
\]
と書く．高分母は $a=sS+1,b=zT+1$，
$x=(a-1)/r,y=(b-1)/(r+1)$ である．
正の枝では
\begin{equation}\label{hs18:positive-bounds}
 S>r\ge1,\quad S\ge2,\quad T\ge S+2,\quad T\ge5,\quad L\ge3.
\end{equation}
中央の負の枝を含む全正局所模型についても $S\ge2,T\ge S$ は成り立つ．
負の枝の $S\ge2$ は式~\eqref{kb15:E-bounds} の $w\ge2$，
$T\ge S$ は補題~\ref{td12:low-entry} による．
以下の $C,D$ は canonical な低 rays の名称であり，出発 core の番号とは別である．

core $3,4$ における ray $2$ は低 ray であり，それぞれ $C$ または $D$ である．
共有高 ray の分母を $q=q_{34}$ とすると，その ray から puncture $2$ へ向かう
entry と reverse entry は次表で尽くされる．
\begin{center}
\begin{tabular}{cccc}
ray $2$ & 共有高 ray & outgoing entry & reverse entry\\\hline
$C$ & $A$ & $x>0$ & $r>0$\\
$C$ & $B$ & $-y<0$ & $-r-1<0$\\
$D$ & $A$ & $-s<0$ & $-S<0$\\
$D$ & $B$ & $z>0$ & $T>0$
\end{tabular}
\end{center}
負の entries $-y,-s$ は $q-1$ の真約数の負であり，$z$ も真約数である．
また $x\ge2$ であり，$x=q-1$ と $r=1$ は同値である．
これらは式~\eqref{hs18:positive-bounds} と pair products から従う．

\begin{lemma}[共有辺 $34$ に残る低 ray の型]
\label{hs18:shared-types}
正局所 charts，整数 clutching，および非負 nodes を仮定する．
共有高辺 $34$ の puncture $2$ においては，次のいずれかだけが残る．
\begin{enumerate}
\item 両 core で ray $2=C$ であり，通常の異符号接着である．
一方の共有高 ray は $A$，他方は $B$ である．
\item 一方の core $i$ で ray $2=D$，他方 $j$ で ray $2=C$ であり，
通常接着である．共有高 ray は両方で $A$ であり，
\[
 r_j=S_i,\qquad 2\le S_i\le4.
\]
\item 一方の core $i$ で ray $2=D$，他方 $j$ で ray $2=C$ であり，
正の例外接着である．共有高 ray は $i$ 側で $B$，$j$ 側で $A$ であり，
\[
 z_i=1,\quad T_i=q-1,\quad r_j=1,\quad x_j=q-1.
\]
\end{enumerate}
両 core で ray $2=D$ となる場合はない．
\end{lemma}

\begin{proof}
共通低 puncture の二つの outgoing entries は $q-1$ の signed divisors である．
整数 clutching により，異符号なら和が零となり，同符号なら絶対値の組は
$\{1,q-1\}$ となる（補題~\ref{ls11:low-puncture}）．
負の entries はいずれも真約数なので，負の例外は存在しない．

両 ray $2=C$ なら，正の entries はともに $x\ge2$ なので正例外も存在せず，第一の場合となる．
両 ray $2=D$ なら，正の entries はともに真約数 $z$ なので正例外はない．
異符号接着では一方の reverse が $T\ge5$ となり，
共通低 puncture の通常 reverse 上界四（補題~\ref{cl16:ordinary-bound}）に反する．

$C,D$ が一つずつの場合，通常接着で $C$ 側が $B$，$D$ 側が $B$ だと，
やはり $T\ge5$ が上界四に反する．従って通常接着では両共有 rays が $A$ となり，
reverse の相殺から $r_j=S_i\le4$ を得る．
正例外では $C$ 側が $A$，$D$ 側が $B$ である．
$z_i$ は真約数なので，最大約数は $x_j$ に限られ，第三の場合を得る．
\end{proof}

\subsection{node 条件による唯一の枝への帰着}

graph の自己同型 $(0\ 1)(3\ 4)$ により，中央 core $2$ の canonical 高 rays を
$A=0,B=1$ としてよい．低 rays は $(C,D)=(3,4)$ または $(4,3)$ である．
その parameters を添字 $2$ で表す．低辺に沿う node 式は
\begin{equation}\label{hs18:low-node}
 n_{2i;k}=1-p^2_{ik}-p^i_{2k}\in\{0,1,2,3,4,5\}.
\end{equation}
特に canonical $D$ row の低低 entry は $-L$ であり，零と置いてはいけない．

\begin{lemma}[node 条件を通過する唯一の中央族]
\label{hs18:node-family}
前補題の仮定の下で，$G_{5,5}$ が接着可能なら，標識交換後に次が必要である：
\[
 \begin{array}{c|c|ccccc}
 \text{core}&(A,B,C,D)&r&s&z&S&T\\\hline
 2&(0,1,4,3)&1&2t-1&t&2&6\\
 3&(0,4,1,2)&1&1&1&2&6
 \end{array}\qquad (t\ge1).
\]
従って
\[
 (q_{02},q_{12},q_{03},q_{34})=(4t-1,6t+1,3,7).
\]
core $4$ の canonical 標識は $(A,B,C,D)=(3,1,2,0)$ であり，
その $a=q_{34}=7,r=1$ が固定される．
\end{lemma}

\begin{proof}
補題~\ref{hs18:shared-types} の三場合を順に調べる．

\paragraph{両外側で ray $2=C$ の場合．}
このとき $p^3_{21}=p^4_{20}=0$ である．
中央が $(C,D)=(4,3)$ なら $n_{23;1}=1-T_2<0$ となる．
中央が $(C,D)=(3,4)$ の場合，core $3$ の共有 ray が $A$ なら，
通常接着から $r_3=r_4+1\ge2$ なので
$n_{23;4}=1-r_3<0$ となる．共有 ray が $B$ なら，
core $3$ のもう一つの高 ray $0$ は $A$ なので
$n_{23;0}=1-r_2-r_3<0$ となる．従って第一の場合は存在しない．

\paragraph{$C,D$ 間の通常接着．}
$D$ 側を $i$，$C$ 側を $j$ とすると，reverse はそれぞれ $-S_i,S_i$ である．
$S_i\ge2$ なので，中央の低 $C$ が $j$ であると
$p^2_{ji}=0$ から $n_{2j;i}=1-S_i<0$ となる．
従って中央の低 $C$ は $i$，低 $D$ は $j$ である．

$i=3,j=4$ なら，core $3$ の共有 ray $4$ が $A$，もう一つの高 ray $0$ が $B$ なので
$n_{23;0}=1-r_2-T_3<0$ となる．
$i=4,j=3$ なら，core $3$ の共有 ray $4$ が $A$，もう一つの高 ray $0$ が $B$ であり，
\[
 p^2_{30}=-S_2,\qquad p^3_{20}=-(S_4+1).
\]
従って
$n_{23;0}=S_2+S_4+2\ge6$ となって上界五に反する．
これで第二の場合も排除された．

\paragraph{$C,D$ 間の正例外接着．}
$D$ 側を $i$，$C$ 側を $j$ とすると，reverse はそれぞれ $T_i\ge5,1$ である．
中央の低 $C$ が $i$ だと $n_{2i;j}=1-T_i<0$ なので，
中央の低 $D$ が $i$，低 $C$ が $j$ である．

$i=4,j=3$ なら，中央の $p^2_{40}=-S_2$ と
core $4$ の $p^4_{20}=-L_4$ により
$n_{24;0}=1+S_2+L_4\ge6$ となる．
従って $i=3,j=4$ であり，canonical 標識は
\[
 (A,B,C,D)_2=(0,1,4,3),\qquad
 (A,B,C,D)_3=(0,4,1,2),\qquad
 (A,B,C,D)_4=(3,1,2,0)
\]
に限られる．ここでは
\[
 n_{23;0}=1+S_2+S_3\le5
\]
なので $S_2=S_3=2$ である．残る二つの node 数は
\[
 n_{23;1}=T_3-T_2,\qquad n_{23;4}=T_2-T_3
\]
となるから $T_2=T_3$ である．

外側の正の枝から $r_3=1$，例外接着から $z_3=1$ である．
$S_3=2$ を高高 pair equation に代入すると
\[
 T_3=8-\frac4{s_3+1},\qquad s_3\ge1,\qquad 2\mid T_3.
\]
$s_3+1$ は四の正約数で二以上だから，唯一の可能性は
$s_3=1,T_3=6$ である．中央も $S_2=2,T_2=6$ なので負の枝ではなく，
$r_2=1$ である．中央の高高 pair equation は $s_2=2z_2-1$ となる．
$t=z_2$ と置くと表示した族を得る．共有分母は $q_{34}=z_3T_3+1=7$ であり，
core $4$ は正例外の最大約数端なので $r_4=1$ である．
\end{proof}

\subsection{二本の低辺の係数相殺による終結}

\begin{lemma}[固定した $a=7,r=1$ の完全な七模型]
\label{hs18:fixed-seven}
二高分母の正の canonical 整数模型で $a=7,r=1$ とする．
全ての可能性は次の七つであり，いずれも $D_C>0$ である．
\begin{center}
\begin{tabular}{rrrrr|r}
$b$&$s$&$z$&$S$&$T$&$D_C$\\\hline
13&3&2&2&6&4\\
19&1&1&6&18&15\\
19&2&2&3&9&9\\
31&1&2&6&15&18\\
41&2&5&3&8&19\\
43&1&3&6&14&23\\
79&1&6&6&13&40
\end{tabular}
\end{center}
これは相手分母 $b$ に上界を置く分類ではない．
\end{lemma}

\begin{proof}
$s\mid6$ なので $s=1,2,3,6$ であり，$x=6,S=6/s$ である．
pair-product の因数分解は
\[
 z\{2(13-s)-(s+1)T\}=2(s-7),\qquad b=zT+1
\]
である．$s=1,2,3,6$ の場合はそれぞれ
\[
 z(T-12)=6,\qquad z(3T-22)=10,\qquad
 z(T-5)=2,\qquad 7z(T-2)=2.
\]
最後は不可能である．前の三つについて，右辺の全正約数 $z$ を取り，
$T$ の整数性と $2\mid zT$ を課すと，表示した七模型をちょうど得る．
それぞれ元の row sums，pair products，正値性を満たす．
$D_C=T-S-x+y$，$y=zT/2$ から最後の列を計算する．
\end{proof}

\begin{theorem}[$G_{5,5}$ support の不存在]
\label{hs18:no-house}
五次の標準 $K_5$ branch graph 上で，正の局所整数 charts，marked puncture の
整数 clutching，非負 nodes，および対向 $D$ の相殺を仮定する．
高分母 support は $G_{5,5}=\{01,02,03,12,14,34\}$ ではあり得ない．
証明は分母の上界も三高分母の一般分類も使わず，$D$ の相殺は低辺 $23,24$ だけに用いる．
\end{theorem}

\begin{proof}
補題~\ref{hs18:node-family} の唯一の族について，低 rays の係数は
\[
 D^2_3=s_2-z_2-1=t-2,\qquad D^2_4=6-t,\qquad
 D^3_2=s_3-z_3-1=-1
\]
である．低辺 $23$ の相殺は $t-2-1=0$，従って $t=3$ を要求する．
低辺 $24$ の相殺は，core $4$ の ray $2=C$ に対し
\[
 D^4_2=-(6-t)=-3
\]
を要求する．一方，core $4$ は $a=7,r=1$ なので，
補題~\ref{hs18:fixed-seven} によって $D^4_2=D_C>0$ である．
これが最後の矛盾である．
\end{proof}

唯一残る族は低三角形 $234$ の node 条件から得られた必要な部分データであり，
未知の三価 cores $0,1$ の整数行列や $q_{01}$ の存在を主張しない．
実際，二本の低辺で既に矛盾するので，それらの未知データを復元する必要はない．
本定理は現在の標準 cones と branch graph を保つ算術的接着の障害である．
一般の Lagrangian torus fibrations に対する不存在定理には拡張していない．
